\section{Metric Observation Layer}

The metric observation layer is the next executable bridge between source
records and Alignment Scores. It is deliberately vertical-neutral: an
observation identifies the entity, metric, raw value, normalization bounds,
source, citation, date, jurisdiction, sample size, and confidence. The
vertical registry supplies the metric direction, weight, and minimum sample
requirements.

\begin{lstlisting}
raw source record
  -> metric observation with provenance
  -> bounded normalization (inverse metrics are reversed)
  -> sample sufficiency check
  -> confidence-weighted metric aggregate
  -> registry-weighted dimension value
  -> universal Alignment Score contract
\end{lstlisting}

The same machinery can process an NCRA observation such as \texttt{LMCR} and
a corrections observation such as \texttt{DPA}; only the vertical registry,
source mapping, and evidence interpretation differ. A first use case does not
become a system-wide assumption.

Observations that do not meet a metric's minimum sample requirement are
retained for audit and shown as insufficient, but are excluded from the
metric aggregate. Every observation retains source provenance so that a score
can be reproduced from the underlying record set.

This layer does not yet fetch live sources or persist snapshots. Those remain
separate ingestion and storage concerns and must be completed before an
empirical profile is published.
